\documentclass[11pt]{article}
\usepackage{explorer}
\newcommand{\ExplorerLanguage}{FR}
\newcommand{\ExplorerNumber}{04}
\newcommand{\ExplorerKind}{Leçon et méthode}
\begin{document}
\ExplorerTitle{Mathématiques expertes · approfondissement}{Bézout et équations entières}
\ExplorerSection{Avant de commencer}
Algorithme d’Euclide et divisibilité dans \(\mathbb{Z}\).\par
\ExplorerSection{Idées et méthode}
Pour deux entiers positifs a et b, l’algorithme d’Euclide calcule \(d = PGCD(a,b)\). En remontant les divisions, on exprime d comme au + bv avec u et v entiers : c’est une identité de Bézout.\par
L’équation ax + by = c possède des solutions entières si et seulement si d divise c. La nécessité vient du fait que d divise ax + by ; la suffisance vient d’une identité de Bézout multipliée par \(\frac{c}{d}\).\par
Si \((x_{0},y_{0})\) est une solution, toutes les solutions sont \(x = x_{0} + (\frac{b}{d})t\) et \(y = y_{0} - (\frac{a}{d})t\), avec t entier. Les contraintes \(x \geq  0\) et \(y \geq  0\) doivent être traitées ensuite.\par
\ExplorerSection{Exemple expliqué}
Pour \(15x + 21y = 6\), on \(a 21 = 15 + 6\) et \(15 = 2 \times  6 + 3\). Ainsi \(3 = 3 \times  15 - 2 \times  21\). Une solution est \((6,-4)\), et toutes les solutions sont \(x = 6 + 7t, y = -4 - 5t\).\par
\ExplorerSection{À explorer}
Tester plusieurs valeurs de t dans l’exemple, puis vérifier l’équation sans arrondir.\par
\end{document}
